\documentclass[12pt, aspectratio=169]{beamer}

\usetheme[progressbar=frametitle, block=fill]{metropolis}

% --- Paquetes ---
\usepackage[utf8]{inputenc}
\usepackage[spanish]{babel}
\usepackage{amsmath, amssymb}
\usepackage{tikz}
\usetikzlibrary{positioning, arrows.meta, shapes.geometric, backgrounds,
                fit, calc, decorations.pathmorphing}
\usepackage{pgfplots}
\pgfplotsset{compat=1.18}
\usepackage{booktabs}
\usepackage{colortbl}
\usepackage{array}
\usepackage{xcolor}

% --- Paleta de colores ---
\definecolor{alrcc-blue}{RGB}{0, 82, 147}
\definecolor{alrcc-teal}{RGB}{0, 150, 136}
\definecolor{alrcc-orange}{RGB}{210, 90, 15}
\definecolor{alrcc-green}{RGB}{46, 139, 87}
\definecolor{alrcc-red}{RGB}{180, 35, 35}
\definecolor{alrcc-lightblue}{RGB}{220, 235, 250}

\setbeamercolor{palette primary}{bg=alrcc-blue}
\setbeamercolor{progress bar}{fg=alrcc-teal, bg=alrcc-blue!40}
\setbeamercolor{alerted text}{fg=alrcc-orange}
\setbeamercolor{example text}{fg=alrcc-teal}
\setbeamercolor{block title}{bg=alrcc-blue!20, fg=alrcc-blue!80!black}
\setbeamercolor{block body}{bg=alrcc-blue!6}
\setbeamercolor{block title alerted}{bg=alrcc-orange!25, fg=alrcc-orange!80!black}
\setbeamercolor{block body alerted}{bg=alrcc-orange!7}
\setbeamercolor{block title example}{bg=alrcc-green!20, fg=alrcc-green!70!black}
\setbeamercolor{block body example}{bg=alrcc-green!6}

% --- Comandos de la notación ALRCC ---
\newcommand{\qg}[1]{\mathsf{#1}}
\newcommand{\stateket}[3]{\left|#1\right\rangle^{\!#2}_{#3}}
\newcommand{\stategate}[3]{\qg{#1}^{#2}_{#3}}
\newcommand{\ket}[1]{\left|#1\right\rangle}
\newcommand{\bra}[1]{\left\langle#1\right|}
\newcommand{\braket}[2]{\left\langle#1\middle|#2\right\rangle}

% Marcas de check y cruz simples
\newcommand{\yes}{\textcolor{alrcc-green}{$\checkmark$}}
\newcommand{\no}{\textcolor{alrcc-red}{$\times$}}
\newcommand{\maybe}{\textcolor{alrcc-orange}{$\sim$}}

% --- Metadatos ---
\title{ALRCC: Lenguaje algebraico para la representación de circuitos cuánticos}
\subtitle{Motivación \textbullet\ Objetivos \textbullet\ Línea de trabajo \textbullet\ Perspectivas}
\author{Francisco Costa Cano}
\institute{Universidad Internacional de La Rioja (UNIR)}
\date{Mayo 2026}

% ====================================================================
\begin{document}
% ====================================================================

% --- Portada ---
{
\setbeamertemplate{footline}{}
\begin{frame}[noframenumbering]
  \titlepage
\end{frame}
}

% --- Índice ---
\begin{frame}{Agenda}
  \setbeamertemplate{section in toc}[sections numbered]
  \tableofcontents[hideallsubsections]
\end{frame}

% ====================================================================
\section{Motivación}
% ====================================================================

\begin{frame}{El paradigma gráfico domina la computación cuántica}
  \begin{columns}[T]
    \column{0.50\textwidth}
    \textbf{Circuito gráfico estándar}

    \medskip
    \begin{tikzpicture}[scale=0.85,
        wire/.style={draw, thick},
        gate1/.style={rectangle, draw=alrcc-blue, thick, fill=alrcc-lightblue,
            minimum width=0.7cm, minimum height=0.55cm, font=\small},
        ctrl/.style={circle, fill=alrcc-blue, inner sep=2pt},
        tgt/.style={circle, draw=alrcc-blue, thick, inner sep=4pt}
      ]
      % Lineas
      \draw[wire] (0,1.2) -- (5.5,1.2) node[right,font=\scriptsize]{};
      \draw[wire] (0,0)   -- (5.5,0)   node[right,font=\scriptsize]{};
      % Etiquetas
      \node[left, font=\small] at (0,1.2) {$\ket{0}$};
      \node[left, font=\small] at (0,0)   {$\ket{0}$};
      % H sobre qubit 1
      \node[gate1] at (1.0,1.2) {$H$};
      % CNOT
      \node[ctrl] at (2.5,1.2) {};
      \node[tgt]  at (2.5,0)   {};
      \draw[thick, alrcc-blue] (2.5,1.2) -- (2.5,0.18);
      \draw[thick, alrcc-blue] (2.5,-0.18) -- (2.5,-0.4);
      \node[font=\scriptsize, alrcc-blue] at (2.5,0) {$+$};
      % Mediciones
      \draw[thick] (4.2,1.2) -- (4.8,1.2);
      \draw[thick] (4.2,0)   -- (4.8,0);
      \node[rectangle, draw, thick, font=\footnotesize,
        minimum width=0.6cm, minimum height=0.45cm] at (4.5,1.2) {M};
      \node[rectangle, draw, thick, font=\footnotesize,
        minimum width=0.6cm, minimum height=0.45cm] at (4.5,0) {M};
    \end{tikzpicture}

    \medskip
    \begin{itemize}\small
      \item[\yes] Intuitivo para circuitos pequeños
      \item[\yes] Universal en la literatura
      \item[\no] \alert{No escala} para $n > 20$ qubits
      \item[\no] \alert{No se manipula algebraicamente}
    \end{itemize}

    \column{0.46\textwidth}
    \begin{block}{Computación clásica}
      \small
      \textbf{Dos} representaciones equivalentes:
      \begin{itemize}\small
        \item Diagrama de circuito lógico
        \item Expresión booleana permiten \textbf{simplificación algebraica}.
      \end{itemize}
    \end{block}

    \medskip
    \begin{alertblock}{La brecha cuántica}
      \small
      \textbf{¿Por qué no existe el equivalente algebraico para circuitos cuánticos?}
    \end{alertblock}
  \end{columns}
\end{frame}

% ---

\begin{frame}{El problema escala exponencialmente}
  \begin{columns}[T]
    \column{0.48\textwidth}
    \textbf{Crecimiento del número de puertas}

    \medskip
    \begin{tikzpicture}[scale=0.9]
      \begin{axis}[
          xlabel={Número de qubits},
          ylabel={Puertas (escala log)},
          ymode=log,
          xmin=5, xmax=40,
          grid=major,
          legend pos=north west,
          width=\textwidth, height=5.2cm,
          tick label style={font=\tiny},
          label style={font=\small},
          legend style={font=\tiny},
        ]
        \addplot[alrcc-blue, very thick, domain=5:40, samples=40]
        {2^(x/5)};
        \addlegendentry{QFT}
        \addplot[alrcc-orange, very thick, dashed, domain=5:40, samples=40]
        {x^2};
        \addlegendentry{Grover}
        \addplot[alrcc-teal, very thick, dotted, domain=5:40, samples=40]
        {x^3/5};
        \addlegendentry{Shor (modular)}
      \end{axis}
    \end{tikzpicture}

    \column{0.48\textwidth}
    \begin{alertblock}{Problema para $n > 20$ qubits}
      \small
      \begin{itemize}
        \item Gráfico no cabe en una página
        \item La optimización se hace heurística o por fuerza bruta
      \end{itemize}
    \end{alertblock}

    \medskip
    \begin{exampleblock}{Lo que ALRCC proporciona}
      \small
      Una \textbf{expresión algebraica compacta}:
      \[
        C = \stategate{X}{\Gamma_1}{\Lambda_1}
        \cdots
        \stategate{X}{\Gamma_k}{\Lambda_k}
      \]
      manipulable simbólicamente,
      \textbf{independiente de} $n$.
    \end{exampleblock}
  \end{columns}
\end{frame}

\section{Estado de la técnica}

\begin{frame}{Panorama de alternativas existentes}
  \centering
  \small
  \begin{tabular}{lccccc}
    \toprule
    \textbf{Formalismo}     & \textbf{Texto}  & \textbf{simbólico}
                            & \textbf{Formal} & \textbf{Reversibles}                   \\
    \midrule
    Circuito gráfico        & \no             & \no                  & \yes   & \yes   \\
    ZX-calculus             & \no             & \yes                 & \yes   & \yes   \\
    OpenQASM~3              & \yes            & \no                  & \yes   & \yes   \\
    Quipper                 & \no             & \maybe               & \yes   & \yes   \\
    Lambda-cálculo cuántico & \no             & \yes                 & \yes   & \maybe \\
    Hutsell--Greenwood      & \yes            & \maybe               & \no    & \yes   \\
    \midrule
    \textbf{ALRCC}          & \yes            & \yes                 & \maybe & \yes   \\
    \bottomrule
  \end{tabular}
\end{frame}

\section{La propuesta ALRCC}

\begin{frame}{Las bases del lenguaje ALRCC}
  \begin{columns}[T]
    \column{0.33\textwidth}
    \begin{block}{1. Estado indexado}
      \centering
      \[
        \stateket{a}{k}{\gamma}
      \]
      \small
      Ket $\ket{a}$ con el qubit $k$ fijado a $\gamma \in \{0,1\}$.

      \smallskip
      \textit{Ejemplo:}
      \[
        \stateket{a_1 a_2 a_3}{2}{0} = \ket{a_1\,0\,a_3}
      \]
    \end{block}

    \column{0.33\textwidth}
    \begin{block}{2. Puerta localizada}
      \centering
      \[
        \qg{A}_k = I^{\otimes k-1} \otimes \qg{A} \otimes I^{\otimes n-k}
      \]
      \small
      La puerta $\qg{A}$ sobre el qubit $k$ únicamente.

      \smallskip
      \textit{Equivalencia:}
      \[
        \qg{A}_k = \stateket{a}{k}{\qg{A}}
      \]
    \end{block}

    \column{0.33\textwidth}
    \begin{block}{3. Conjuntista}
      \centering
      \[
        \qg{A}_\Gamma = \prod_{k \in \Gamma} \qg{A}_k
      \]
      \small
      $\qg{A}$ aplicada a cada qubit del conjunto $\Gamma \subseteq \mathbb{N}_n$.

      \smallskip
      Bien definida porque puertas en qubits distintos conmutan.
    \end{block}
  \end{columns}
\end{frame}

\begin{frame}{Las bases del lenguaje ALRCC}
  \begin{block}{Extensión: puerta controlada general $\stategate{A}{\Gamma}{\Lambda}$}
    \small
    Con controles $\Gamma$ y targets $\Lambda$, $\Gamma \cap \Lambda = \emptyset$:
    \[
      \stategate{A}{\Gamma}{\Lambda}\ket{a}
      = \ket{a}
      - \Bigl(\prod_{h \in \Gamma} a_{h1}\Bigr)\,\stateket{a}{\Gamma}{\mathbf{1}}
      + \Bigl(\prod_{h \in \Gamma} a_{h1}\Bigr)\,\stateket{a}{\Gamma \cup \Lambda}{\mathbf{1}\,\qg{A}}
      \quad\quad \stategate{A}{\emptyset}{\Lambda} = \qg{A}_\Lambda
    \]
  \end{block}
\end{frame}

% ---

\begin{frame}{El resultado algebraico central}
  \begin{columns}[T]
    \column{0.52\textwidth}
    \begin{block}{Proposición}
      Sea $\qg{A}$ una puerta de 1~qubit y $\Gamma, \Lambda \subseteq \mathbb{N}_n$. Entonces:
      \[
        \boxed{\qg{A}_\Gamma\,\qg{A}_\Lambda
          = \qg{A}_{\Gamma \cup \Lambda}\;\qg{A}_{\Gamma \cap \Lambda}}
      \]
    \end{block}

    \medskip
    \textbf{Corolario para puertas involutivas} ($\qg{A}^2 = \qg{I}$,
    ej.\ $\qg{X}$, $\qg{H}$, $\qg{Z}$):
    \[
      \qg{A}_\Gamma\,\qg{A}_\Lambda = \qg{A}_{\Gamma \triangle \Lambda}
    \]
    donde $\triangle$ es la diferencia simétrica.

    \medskip
    \alert{Dos puertas sobre conjuntos solapados se reducen a una sola.}
    Aplicable recursivamente en cadenas de puertas.

    \column{0.44\textwidth}
    \begin{exampleblock}{Ejemplo de simplificación}
      \small
      Circuito original:
      \[
        \qg{X}_{\{1,2,3\}} \cdot \qg{X}_{\{2,3,4\}}
      \]
      Aplicando resultado
      \[
        \qg{X}_{\{1,4\}}
      \]
    \end{exampleblock}

    \begin{exampleblock}{}
      \centering\small
      \textbf{6 aplicaciones de puerta} $\rightarrow$ \textbf{2 aplicaciones}
    \end{exampleblock}
  \end{columns}
\end{frame}

% ---

\begin{frame}{Demostración de potencia: sumador binario}
  \begin{columns}[T]
    \column{0.54\textwidth}
    \textbf{Circuito que computa $2 + 3 = 5$}

    \textbf{10 qubits}

    \medskip
    En notación ALRCC, el circuito completo en \textbf{una línea}:
    \begin{align*}
      C & = \stategate{X}{8}{10}\,\stategate{X}{4,7}{10}
      \stategate{X}{7}{9}\,\stategate{X}{4}{9}
      \stategate{X}{5,6}{8}\,\stategate{X}{1,2}{4}
      \stategate{X}{6}{7}\,\stategate{X}{2}{3}
      \stategate{X}{5}{7}\,\stategate{X}{1}{3}
    \end{align*}

    Se aplica sobre $\ket{0100110000}$ y la medición de los
    qubits $\{3,9,10\}$ produce $101$.

    \column{0.42\textwidth}
    \textbf{Comparación de tamaño de descripción}

    \begin{tikzpicture}[scale=0.88]
      \begin{axis}[
          xlabel={Bits sumados ($n$)},
          ylabel={Líneas de descripción},
          xmin=1, xmax=8,
          ymin=0, ymax=500,
          grid=major,
          legend pos=north west,
          width=\textwidth, height=5.5cm,
          tick label style={font=\tiny},
          label style={font=\small},
          legend style={font=\tiny},
        ]
        \addplot[alrcc-blue, very thick, mark=*]
        coordinates {(1,2)(2,4)(3,6)(4,8)(5,10)(6,12)(7,14)(8,16)};
        \addlegendentry{ALRCC ($O(n)$)}
        \addplot[alrcc-red, very thick, dashed, mark=square*]
        coordinates {(1,4)(2,20)(3,55)(4,120)(5,220)(6,360)(7,450)(8,500)};
        \addlegendentry{Circuito gráfico ($O(n^2)$)}
      \end{axis}
    \end{tikzpicture}
  \end{columns}
  \begin{alertblock}{Reto actual}
    \small
    La demostración general para $n$~bits requiere las
    reglas de reescritura completas del cálculo.
    Es uno de los objetivos centrales del trabajo.
  \end{alertblock}
\end{frame}

% ====================================================================
\section{Objetivos de investigación}
% ====================================================================

\begin{frame}{Objetivo principal y sub-objetivos}
  \begin{block}{Objetivo principal}
    \small
    \textbf{Cálculo algebraico formal} para circuitos cuánticos
    reversibles que permita expresar, comparar y \textbf{optimizar}
    circuitos de $n$~qubits mediante reescritura simbólica.
  \end{block}

  \medskip
  \begin{columns}[T]
    \column{0.48\textwidth}
    \textbf{Sub-objetivos técnicos}
    \begin{enumerate}\small
      \item[\textbf{O1}] Sintaxis formal (BNF).
      \item[\textbf{O2}] Sistema de reglas de reescritura con terminación
            y confluencia.
      \item[\textbf{O3}] Demostración de \emph{soundness} respecto a la
            semántica matricial.
      \item[\textbf{O4}] Optimización para circuitos Toffoli.
      \item[\textbf{O5}] Entrelazamiento y mediciones.
    \end{enumerate}

    \column{0.48\textwidth}
    \textbf{Sub-objetivos aplicados}
    \begin{enumerate}\small
      \item[\textbf{O6}] Prototipado a Mathematica.
      \item[\textbf{O7}] Aplicación a códigos de correción cuántica de errores.
      \item[\textbf{O8}] Aplicación algoritmos criptográficos cuánticos.
    \end{enumerate}
  \end{columns}
\end{frame}


% ====================================================================
\section{Conclusiones}
% ====================================================================

\begin{frame}{Resumen}
  \begin{columns}[T]
    \column{0.48\textwidth}
    \textbf{Estado actual del trabajo}
    \begin{itemize}\small
      \item[\yes] Notación operativa y demostrada en casos concretos
      \item[\yes] Identidad de conjuntos: herramienta de optimización real
      \item[\yes] Sumador binario representado en una sola expresión
      \item[\yes] Conmutación, involución y eigenvalores demostrados
      \item[\yes] Análisis crítico de posicionamiento realizado
    \end{itemize}

    \column{0.48\textwidth}
    \textbf{Próximos pasos}
    \begin{itemize}\small
      \item[\textbf{1.}] Formalizar BNF + semántica
      \item[\textbf{2.}] Demostración general del sumador $n$-bits
    \end{itemize}

    \begin{block}{Mensaje central}
      \small
      ALRCC es una propuesta \textbf{original, necesaria y sólida}: lenguaje algebraico textual para circuitos cuánticos
      clásico-reversibles, \textbf{complementario a ZX-calculus}.
    \end{block}
  \end{columns}
\end{frame}

{
\setbeamertemplate{footline}{}
\begin{frame}[noframenumbering]
  \centering
  \vspace{1.5cm}
  {\Large \textbf{Gracias por su atención}}

  \vspace{0.8cm}
  {\large Francisco Costa Cano}\\[0.3cm]
  {\normalsize francisco.costacano@unir.net}
\end{frame}
}

\end{document}
